% META: {"id": "TEXP-ARI-EX-010", "chapitre": "TEXP-ARITHMETIQUE", "type_objet": "exercice", "capacites": ["TEXP-ARI-C2"], "capacites_codes": ["C2"], "methodes": ["M2"], "parcours": 1, "competences": ["calculer"], "duree_min": 12, "mode_creation": "ex_nihilo", "sources_inspiration": [], "fichier_tex": "chapitres/TEXP-ARITHMETIQUE/exercices/TEXP-ARI-EX-010.tex", "corrige_tex": "chapitres/TEXP-ARITHMETIQUE/corriges/TEXP-ARI-CO-010.tex", "status": "generated"}
% BEGIN-VERIFY
% table = [(5 * k) % 12 for k in range(12)]
% assert table == [0, 5, 10, 3, 8, 1, 6, 11, 4, 9, 2, 7]
% assert (5 * 5) % 12 == 1
% assert [x for x in range(12) if (5 * x) % 12 == 3] == [3]
% assert (5 * (3 + 12 * 4)) % 12 == 3
% END-VERIFY
\begin{exercice}{TEXP-ARI-EX-010}{1}{12}
\begin{enumerate}
  \item Dresser la table des restes de $5k$ modulo $12$ pour
        $k\in\{0,1,\dots,11\}$.
  \item En déduire l'inverse de $5$ modulo $12$.
  \item Résoudre $5x\equiv3\ [12]$.
\end{enumerate}
\end{exercice}
